% SEGMENT OLP-0472-S01
% Part: normal-modal-logic
% Chapter: sequent-calculus
% Section: rules-for-K

\documentclass[../../../include/open-logic-section]{subfiles}

\begin{document}

\olfileid{nml}{seq}{rul}

\olsection{\Log{K} க்கான விதிகள்}

வழக்கமான கூற்றுத் தருக்க இணைப்புகளுக்கான விதிகள், வழக்கமான தொடரணிக்
கணியம்~$\Log{LK}$ இல் உள்ளவற்றைப் போன்றவையே. அடிகோள்களும் அவ்வாறே:
$!A \Sequent !A$ என்ற வடிவிலுள்ள ஒவ்வொரு தொடரணியும் ஓர் அடிகோளாகக்
கொள்ளப்படுகிறது.

வாய்ப்புநிலைச் செயலி\iftag{prvBox}{\iftag{prvDiamond}{களான~$\Box$
மற்றும்}{யான~$\Box$}}{}%
\iftag{prvDiamond}{யான~$\Diamond$}{} க்கு, பின்வரும் கூடுதல்
\iftag{notprvBox,notprvDiamond}{விதி}{விதிகள்} உள்ளன:
\[
  \iftag{prvBox}{\iftag{prvDiamond}{% <> and [] primitive
    \Axiom$\Gamma \fCenter \Delta, !A$
    \RightLabel{$\Box$}
    \UnaryInf$\Box\Gamma \fCenter \Diamond\Delta, \Box!A$
    \DisplayProof
    \qquad
    \Axiom$!A, \Gamma \fCenter \Delta$
    \RightLabel{$\Diamond$}
    \UnaryInf$\Diamond!A, \Box\Gamma \fCenter \Diamond\Delta$
    \DisplayProof
}{% only Box primitive
\Axiom$\Gamma \fCenter !A$
\RightLabel{$\Box$}
\UnaryInf$\Box\Gamma \fCenter \Box!A$
\DisplayProof
}}{% only <> primitive
  \Axiom$!A \fCenter \Delta$
  \RightLabel{$\Diamond$}
  \UnaryInf$\Diamond!A \fCenter \Diamond\Delta$
  \DisplayProof
}
\]
இங்கே, \iftag{prvBox}{$\Box\Gamma$ என்பது $\Gamma$ விலுள்ள ஒவ்வொரு
!!{formula}க்கும் முன் $\Box$ ஐ இடுவதால் $\Gamma$ விலிருந்து கிடைக்கும்
!!{formula}s இன் தொடர்முறை\iftag{prvDiamond}{; மேலும்
}{}}{}%
\iftag{prvDiamond}{$\Diamond\Delta$ என்பது $\Delta$ விலுள்ள ஒவ்வொரு
!!{formula}க்கும் முன் $\Diamond$ ஐ இடுவதால் $\Delta$ விலிருந்து கிடைக்கும்
!!{formula}s இன் தொடர்முறை}{}.
\iftag{notprvDiamond}{$\Box$ விதியின் மேல்த் தொடரணியின் வலப்பக்கத்தில்
அதிகபட்சமாக ஒரே ஓர் !!{formula}~$!A$ மட்டுமே இருக்க வேண்டும். }{}%
\iftag{notprvBox}{$\Diamond$ விதியின் மேல்த் தொடரணியின் இடப்பக்கத்தில்
அதிகபட்சமாக ஒரே ஓர் !!{formula}~$!A$ மட்டுமே இருக்க வேண்டும். }{}%
\iftag{prvBox}{$\Gamma$\iftag{prvDiamond}{ வும்~}{}}{}%
\iftag{prvDiamond}{$\Delta$ வும்}{} வெற்றாக இருக்கலாம்; அப்போது கீழ்த்
தொடரணியின் அதற்குரிய
\iftag{prvBox}{$\Box\Gamma$\iftag{prvDiamond}{ பகுதியும்~}{}}{}%
\iftag{prvDiamond}{$\Diamond\Delta$ பகுதியும்}{} வெற்றாகவே இருக்கும்.

ஒரே ஓர் !!{formula}~$!A$ க்கு மட்டும்
\iftag{prvBox}{வலப்பக்கத்தில் $\Box$ ஐ\iftag{prvDiamond}{யும்
}{}}{}%
\iftag{prvDiamond}{இடப்பக்கத்தில் $\Diamond$ ஐயும்}{} சேர்க்க வேண்டும்
என்ற கட்டுப்பாடு அவசியமானது.
\iftag{prvBox}{வலப்பக்கத்தில் எத்தனை !!{formula}s க்கு வேண்டுமானாலும்
$\Box$ ஐச் சேர்க்க\iftag{prvDiamond}{வோ, }{}}{}%
\iftag{prvDiamond}{இடப்பக்கத்தில் எத்தனை !!{formula}s க்கு வேண்டுமானாலும்
$\Diamond$ ஐச் சேர்க்கவோ}{} அனுமதித்தால், பின்வருவதை
!!{derive} முடியும்:
  \[
    \iftag{prvBox}{
      \Axiom$!A \fCenter !A$
      \RightLabel{\RightR{\lnot}}
      \UnaryInf$\fCenter !A, \lnot !A$
      \RightLabel{$\Box*$}
      \UnaryInf$\fCenter \Box!A, \Box\lnot !A$
      \doubleLine
      \RightLabel{\RightR{\lor}}
      \UnaryInf$\fCenter \Box!A \lor \Box \lnot !A$
      \DisplayProof}{}
    \iftag{notprvBox,notprvDiamond}{}{\qquad}
    \iftag{prvDiamond}{
      \Axiom$!A \fCenter !A$
      \RightLabel{\LeftR{\lnot}}
      \UnaryInf$\lnot !A, !A \fCenter$
      \RightLabel{$\Diamond*$}
      \UnaryInf$\Diamond\lnot !A,\Diamond!A \fCenter $
      \RightLabel{\RightR{\lnot}}
      \UnaryInf$\Diamond!A \fCenter \lnot\Diamond\lnot !A$
      \RightLabel{\RightR{\lif}}
      \UnaryInf$ \fCenter \Diamond!A \lif \lnot\Diamond\lnot !A$
      \DisplayProof}{}
  \]
ஆனால் \iftag{prvBox}{$\Box!A \lor \Box \lnot !A$
\iftag{prvDiamond}{ மற்றும் $\Diamond!A \lif \lnot\Diamond\lnot !A$
ஆகியவை}{ என்பது}}{$\Diamond!A \lif \lnot\Diamond\lnot !A$ என்பது}
~$\Log{K}$ இல் செல்லுபடியாகாதது.

மேல்த் தொடரணியில் $!A$ உடன் பக்க வாய்பாடுகளையும் அனுமதித்து,
\iftag{prvBox}{$\Box$ விதி வலப்பக்கத்திலுள்ள $!A$ க்கு மட்டும் $\Box$ ஐச்
சேர்க்க\iftag{prvDiamond}{ அனுமதித்தாலோ, }{}}{}%
\iftag{prvDiamond}{$\Diamond$ விதி இடப்பக்கத்திலுள்ள $!A$ க்கு மட்டும்
$\Diamond$ ஐச் சேர்க்க அனுமதித்தாலோ}{} (பக்க வாய்பாடுகளை மாற்றாமல்
விட்டால்), பின்வருவதை !!{derive} முடியும்:
\[
  \iftag{prvBox}{
    \Axiom$!A \fCenter !A$
    \RightLabel{\RightR{\lnot}}
    \UnaryInf$\fCenter !A, \lnot !A$
    \RightLabel{\RightR{\Exchange}}
    \UnaryInf$\fCenter \lnot !A, !A$
    \RightLabel{$\Box*$}
    \UnaryInf$\fCenter \lnot!A, \Box !A$
    \doubleLine
    \RightLabel{\RightR{\lor}}
    \UnaryInf$\fCenter \lnot!A \lor \Box !A$
    \DisplayProof}{}
  \iftag{notprvBox,notprvDiamond}{}{\qquad}
  \iftag{prvDiamond}{
    \Axiom$!A \fCenter !A$
    \RightLabel{\LeftR{\lnot}}
    \UnaryInf$\lnot !A, !A \fCenter$
    \RightLabel{$\Diamond*$}
    \UnaryInf$\Diamond\lnot !A, !A \fCenter $
    \RightLabel{\RightR{\lnot}}
    \UnaryInf$!A \fCenter \lnot\Diamond\lnot !A$
    \RightLabel{\RightR{\lif}}
    \UnaryInf$ \fCenter !A \lif \lnot\Diamond\lnot !A$
    \DisplayProof}{}
\]
ஆனால் \iftag{prvBox}{$\lnot!A \lor \Box !A$ ($!A \lif \Box!A$ க்கு
நிகரானது)\iftag{prvDiamond}{ மற்றும்
$!A \lif \lnot\Diamond\lnot !A$ ஆகியவை}{ என்பது}}%
{$!A \lif \lnot\Diamond\lnot !A$ என்பது}~$\Log{K}$ இல்
செல்லுபடியாகாதது.

\end{document}
